\section*{第 \ref{ch:g9:gravitation} 章 --- 万有引力}

\begin{solution}{exo:g9:gravitation:1}
$F = G \, m_A m_B / d^2$，质量以千克计，$d$ 为
两中心之间的距离（以米计），$F$ 以牛顿计，并且
$G = \qty{6.67e-11}{N.m^2/kg^2}$。用文字来说：引力随每个质量的增大而增大，
随距离的平方增大而减弱——
并且它对两个物体的作用大小相等、方向相反。
\end{solution}

\begin{solution}{exo:g9:gravitation:2}
恰为 \qty{1}{N}——引力定律中的作用力总是成对出现，
大小相等、方向相反。苹果明显加速，是因为
它的质量很小；同样的 \qty{1}{N} 作用在
\qty{5.97e24}{kg} 的行星上，几乎无法察觉其运动变化。
\end{solution}

\begin{solution}{exo:g9:gravitation:3}
在 $2d$ 处：引力为 $d$ 处的四分之一——二的平方的倒数。
在 $10d$ 处：引力为原来的百分之一。
\end{solution}

\begin{solution}{exo:g9:gravitation:4}
$F = \num{6.67e-11} \times 80 \times 80 \div 0.25 =
\qty{1.7e-6}{N}$——不足百万分之二牛顿：
无论把舞伴拉到一起的是什么力，它都不是
万有引力。
\end{solution}

\begin{solution}{exo:g9:gravitation:5}
小球下落的快慢，以及
球形地球表面在其下方弯离的快慢。轨道即
一边下落、一边错过地面的运动。
\end{solution}

\begin{solution}{exo:g9:gravitation:6}
在空间站所在的高度，重力仍保持
接近地面的全部强度——正是它使空间站
维持在圆周轨道上。空间站与宇航员都在\emph{下落}，
共同绕地球下落：一起下落的物体彼此之间不相互挤压，
所谓漂浮，正是共同自由下落时的感觉。
\end{solution}

\begin{solution}{exo:g9:gravitation:7}
因为天体的质量极其巨大——行星与
恒星把微弱的 $G$ 放大为强大的力——
而引力从不被抵消：它只会吸引，因此
每增加一份质量，引力便加深。
\end{solution}

\begin{solution}{exo:g9:gravitation:8}
月球的引力使海水向两侧隆起：朝向月球的一侧（离月球最近的水，
受力最大）与背离月球的一侧（离月球最远的水，
受力最小）；自转的地球每天让每处海岸都经过这两处
隆起——形成两次潮汐。在新月与满月时最大，
此时太阳与月球的引力方向一致：形成大潮。
\end{solution}

\begin{solution}{exo:g9:gravitation:9}
$F = \num{6.67e-11} \times (\num{7.3e22} \times
\num{5.97e24}) \div (\num{3.8e8})^2 \approx \qty{2.0e20}{N}$
——即一个 $2$ 后面跟着二十个零的牛顿，
把月球维系在它的轨道上。
\end{solution}

\begin{solution}{exo:g9:gravitation:10}
两个因素相互反向作用：月球的质量约小八十倍，
使引力减小为八十分之一；但它的
半径约小四倍，却\emph{增大}了表面
引力——增大为四的平方倍，约十四倍。八十倍的减弱，
与十四倍的增强相抵，约剩六分之一。
\end{solution}

\begin{solution}{exo:g9:gravitation:11}
平方反比使太阳的束缚随距离迅速减弱，而较弱的束缚
只能维持较慢的绕行：外侧的
行星被松散地拴住，缓缓运行；内侧的
行星被紧紧抓住，飞快地运行——水星的疾驰与海王星的
164 年的漫步，是同一定律在两种距离下的体现。
\end{solution}

\begin{solution}{exo:g9:gravitation:12}
它的周期必须恰好为一天，并且向东运行——与
地球自转的方向相同。于是卫星与天线完全同步地转动，
“悬停”的碟形天线，实际上是在注视一颗炮弹，
它绕行地球下落，其速率恰好等于地球
在其下方的自转速率——只有在随地球旋转的
观察者看来，它才是静止的。
\end{solution}

\begin{solution}{pb:g9:gravitation:1}
\textbf{1.} $F = \num{6.67e-11} \times (0.1 \times
\num{5.97e24}) \div (\num{6.4e6})^2 \approx \qty{0.97}{N}$
——可称为一牛顿：苹果所受的拉力，
由宇宙常数算得。
\textbf{2.} 同样为 \qty{0.97}{N}——相互性：定律中的
引力总是成对出现，大小相等、方向相反。
\textbf{3.} 中心距离加倍，引力变为四分之一：
约 \qty{0.24}{N}。
\textbf{4.} 行星的每一小块都吸引着苹果——
近处的山脉、深处的地核、远处的对跖点——所有这些
引力之和，恰等同于全部质量集中于
球心时的引力：这是一个定理，而非假设；其作者
直到能够证明它，才推迟发表了多年。
\textbf{5.} 约 $2000 \times 9.8 \approx \qty{2.0e4}{N}$
——即两万牛顿（若取 $d$ 为一个地球
半径，由定律给出的结果相同）。
\textbf{6.} 在两倍半径处，平方反比使引力只剩
四分之一：约 \qty{4900}{N}，相对于 \qty{19600}{N}
的地表重力。
\textbf{7.} 是的——满额重力的四分之一绝非
无重力的真空。卫星把这份引力用于弯曲
自己的轨迹：它绕地球下落，永远错过地面，正如
山顶大炮的例子所示。
\textbf{8.} “足够快地发射，炮弹下落的
轨迹恰与地球表面的弯曲相符，于是绕行不止——
只是下落而不着地。每颗卫星都是这样的炮弹，
横向发射得足够快，从而不断错过地面。”
\textbf{9.} $\num{3.8e8} \div \num{6.4e6} \approx 59.4$
约 59.4 个地球半径。
\textbf{10.} 约为 $1 \div 59.4^2 \approx 1/3500$
的地表强度。
\textbf{11.} 之所以平缓，是因为六十个地球半径处的重力
被稀释为三千五百分之一——五米的
第一秒下落缩短为约一毫米。之所以永不停歇，
是因为这一毫米的下落，恰与
月球横向滑行时地球表面在其下方弯离的距离相等：
下落的距离从不累积。
\textbf{12.} 例如：“苹果、卫星与月球都遵循
同一句话：任意两个质量都向对方下落，
其大小为 $G m_A m_B / d^2$——苹果落地，卫星与
月球则不断错过地面，二者的区别只在于
横向速度。”
\end{solution}
